\documentclass{article}
\usepackage[utf8]{inputenc}
\usepackage{geometry}
\usepackage{graphicx}
\usepackage{booktabs}
\usepackage{hyperref}

\geometry{a4paper, margin=1in}

\title{\textbf{Progress Report}}
\author{Tongyan CONG}
\date{December 12, 2025}

\begin{document}

\maketitle

\section{Clarification: What is ``Background'' in Detection Results?}
In response to the previous discussion, I provide a detailed explanation of the \textbf{background} category observed in the confusion matrix.

\noindent In the context of object detection, \textbf{background} does not refer to a specific visual element in the document, nor is it an additional object class defined by the dataset. Instead, background is an \textit{implicit category} representing all regions where no target object is detected. In other words, any pixel area that does not belong to one of the annotated classes—such as title, author, section header, paragraph, table, figure, formula, or footnote—is automatically treated as background by the model.

\noindent During evaluation, background acts as the destination for \textit{missed detections}. If the model fails to identify an existing object (e.g., a formula), the prediction for that region falls into the background category. This is why a low recall for formulas appears as ``formula $\rightarrow$ background'' in the confusion matrix. Thus, background serves as the model's representation of ``no object here,'' and is essential for analyzing false negatives and understanding detection bottlenecks.




\section{Progress on Reproducing the Baseline System (Sulaeman, 2019)}

\subsection{Reproduction Results and Identified Issues}
The reproduction experiments revealed:
\begin{itemize}
    \item \textbf{Paragraph text extraction is relatively stable}, fully usable for LaTeX assembly.
    \item \textbf{Title and section header recognition still has errors}, especially when fonts are small or formatting is complex.
    \item \textbf{Two-column reading order is the current major bottleneck.} Simple vertical sorting cannot correctly recover reading order, causing paragraph misplacement. This requires a dedicated ``column segmentation'' module as follow-up work.
    \item \textbf{Tables and formulas require further processing} (the original thesis did not fully resolve these either).
\end{itemize}


\begin{figure}[ht]
    \centering
    \begin{minipage}{0.48\textwidth}
        \centering
        \includegraphics[width=\textwidth]{original_dataset-1.png}
        \caption{Original Dataset Document}
    \end{minipage}
    \hfill
    \begin{minipage}{0.48\textwidth}
        \centering
        \includegraphics[width=\textwidth]{reproduced_output-1.png}
        \caption{Reproduced Output}
    \end{minipage}
    \label{fig:comparison}
\end{figure}

\subsection{Insights from the Reproduction Work}
The reproduction process yielded several important insights:

\begin{enumerate}
    \item \textbf{Visual detection is no longer the system bottleneck.} YOLOv8's detection performance is satisfactory. Future focus should shift to text understanding and structural recovery.
    
    \item \textbf{OCR and reading order determination are critical for generation quality.} These modules directly affect whether paper content can be correctly transformed into structured representations.
    
    \item \textbf{Stronger layout understanding capabilities are needed.} Two-column documents, title hierarchies, and paragraph attribution currently rely on simple heuristics, which show limited performance.
\end{enumerate}



\section{Insights from Recent Literature}

\subsection{Paper Overview}
\textbf{Title:} ``Automatic Text Box Placement for Supporting Typographic Design''

\noindent \textbf{Authors:} Jun Muraoka, Daichi Haraguchi, Naoto Inoue, Wataru Shimoda, Kota Yamaguchi, and Seiichi Uchida

\noindent This study investigates the task of predicting optimal text-box coordinates within graphic design layouts. The authors compare task-specific Transformers against fine-tuned and zero-shot Vision-Language Models (VLMs), demonstrating that specialized spatial reasoning models outperform general-purpose multimodal systems in precise layout prediction.

\subsection{Key Implications for My Research}

\subsubsection{Separation of Semantic Understanding and Spatial Layout Modeling}
The paper demonstrates that a task-specific Transformer significantly outperforms both fine-tuned and zero-shot VLMs in precise text-box localization. This suggests that, in my system, LLMs/VLMs should focus on \textit{content understanding}—extracting structure, summarizing sections, and interpreting figures—while a dedicated Transformer-based module should handle \textit{spatial layout prediction}. Attempting to let a single multimodal LLM generate both content and coordinates end-to-end is unlikely to yield stable layouts.

\subsubsection{Element-Level Visual Features Are More Effective Than Holistic Page Images}
The ``multi-image Transformer'' variant achieves noticeably higher accuracy by encoding each layout element (text, image, mask, shape, etc.) as an individual bitmap. This directly informs the design of my slide-generation pipeline: figures, tables, equations, and paragraph blocks from a paper should be \textit{isolated as independent image patches}, paired with textual metadata, and provided to the layout model in a structured JSON format. This ensures the model can reason about inter-element relationships, whitespace availability, and visual saliency more precisely than using only a full-page screenshot.

\subsubsection{Layout Prediction Is Inherently Non-Unique}
The paper highlights that many layouts admit multiple valid placements. Future models may benefit from producing a probability distribution or multiple candidate bounding boxes rather than a single deterministic coordinate. This insight is highly relevant to slide generation, where stylistic flexibility is important. My system should therefore support \textit{multi-solution layout proposals}—e.g., top-k candidate placements or a heatmap of feasible regions—enabling downstream selection or user-guided refinement.

\subsubsection{Practical Data and Model-Input Pipeline Template}
The Crello dataset and its layered representation (e.g., \texttt{textElement}, \texttt{imageElement}, \texttt{svgElement}), together with the JSON schema for element attributes (\texttt{left}, \texttt{top}, \texttt{width}, \texttt{height}), provide a ready-made template for building my own multimodal slide dataset. The paper's input formulation—structured metadata plus per-element visual features—can serve as a strong baseline for constructing the training pipeline in my system.





\end{document}
